%For the simplified model discussed in Section \ref{sec:simple-model}, the mixing term is ignored, and only advection is considered in a simplified manner.

\begin{itemize}
    \item Introduce simplified model (used to understand general trends that can be expected from advection) (Appendix for derivation of non-linear least-squares formatted equation)
    \item Show flow diagram and discuss (several paragraphs to fully explain)
    \item Discuss specific considerations for OpenMC simulation (self-shielding)
    \item Discuss specifics of OpenMC model
    \item Discuss consistent source particle rate between simulations (with equation)
\end{itemize}


\begin{comment}
After the OpenMC model, the custom scripts handle converting the concentrations into group parameters.
Within these scripts, the SciPy non-linear least squares fitting function is used to generate these group parameters.
However, there are some alterations included to ensure the fit is as good as possible.
The first is linearization and normalization of the data.
The normalization is handled by taking the delayed neutron counts, dividing by the fission term \footnote{The fission term is either the net number of fissions or the fission rate, depending on the irradiation type.}, and then taking the natural logarithm.
These normalized counts are then passed into the non-linear least squares fitting.
The linearization is handled by taking the natural logarithm of the entire delayed neutron count fit function.
This treatment reduces both the covariance and the $\chi^2$ of the resulting non-linear least squares fit.
Bounds are also set from zero to 1000 for each term to speed up the process, while the tolerances for convergence are set to $2.23 \times 10^{16}$ to ensure the converged solution is well-refined.

Uncertainties are tracked in this work using the uncertainties package \cite{uncertaintiespackage}.
Uncertainties in the concentrations generated from OpenMC are not accounted for, but uncertainties in the experimentally measured emission probabilities and half lives are.
These uncertainties are used in the minimization of $\chi^2$, shown in Equation \eqref{eq:chisq-nnls}.

\begin{equation}
    \chi^2 = \sum_{t=0}^{t_f} \left(  \frac{b_t - A_t x_t}{\sigma_t}  \right)
    \label{eq:chisq-nnls}
\end{equation}

The terms in Equation \eqref{eq:chisq-nnls} are defined as follows:

\begin{description}
    \item $t_f$ is the final time after irradiation, starting at $t=0$ immediately following irradiation
    \item $b_t$ is the delayed neutron count vector at the row where the time value is $t$
    \item $A_t$ is the coefficient matrix at the row where the time value is $t$
    \item $x_t$ is the solution vector at the row where the time value is $t$
\end{description}

\end{comment}